\Article{22}
The preceding experiments were undertaken chiefly to put the law of magnetic action beyond all doubt. They were also meant to show how many terms of the series should be included and what precision the experiments allow. They showed that if we keep the distances from \SourceLeaf{115}falling below four times the length of the needles, two terms of the series suffice\HistoricalNote{12}. In any case, the differences yielded by the calculation must by no means be treated simply as observational errors. Several precautions from which we may expect still closer agreement were not yet ready at that time. Among them are corrections for hourly variations in the intensity of terrestrial magnetism. These must be taken into account with another needle for comparison, following the method we discussed in Article $10$, II. To ascertain the value of terrestrial magnetism insofar as it can be derived from \textit{these} experiments, we add an account of the remaining experiments that bear on it.

For the first needle and the thread from which it hung, the value of the fraction $\frac{\theta}{Tm}$ was found by the method given in Article $8$ to be $=\frac{1}{251{,}96}$. This gives
\[
\frac MT=0{,}0436074
\]
This number uses the meter as the unit of distance. If we prefer to use the millimeter as the unit, we must multiply this number by the cube of one thousand, giving
\[
\frac MT=43607400
\]

For the second needle, experiments carried out on June $28$, exactly like those we discussed for another needle in Article $11$, gave the following result when the millimeter, the milligram and the second of mean solar time were taken as the units.
\[
TM=135457900
\]
Eliminating the quantity $M$ then gives
\[
T=1{,}7625
\]
